\chapter{Au-delà d'Almgren--Chriss}\label{ch:mx:beyond-almgren-chriss}

Le calendrier d'Almgren--Chriss ne regarde pas le marché pendant qu'il exécute. Un trader qui anticipe une baisse du cours dans la demi-heure à venir doit vendre
plus vite~; celui dont l'impact se résorbe doit marquer une pause après une grosse transaction~; celui dont le marché s'assèche à midi doit négocier avant que cela ne se produise. Ce chapitre ajoute
au bloc du chapitre~14 chacun de ces éléments~: une prévision, les mouvements propres du cours, l'impact transitoire et la liquidité aléatoire, et mesure ce que vaut chacun.

\section{Exécution avec un signal}

\begin{definition}[Calendrier adapté au signal]\label{def:mx:beyond-almgren-chriss:signal}
Un \emph{calendrier adapté au signal}\index{calendrier adapté au signal} (signal-adaptive schedule) fixe la transaction de chaque intervalle à partir des quantités restant à exécuter et d'une prévision de la dérive du cours,
en allant plus vite lorsque la prévision est contraire à l'ordre (un vendeur qui anticipe une baisse) et plus lentement lorsqu'elle va dans son sens.
\end{definition}

Supposons que la dérive du cours, de taux $\alpha_t$, suive un processus d'Ornstein--Uhlenbeck (le livre~4, chapitre~4), avec une demi-vie et un
écart-type stationnaire, observé par le trader. Sur $n$ pas de longueur $\tau$, en vendant $u_k$ à partir des quantités $x_k$ détenues, le coût espéré est quadratique dans l'état
$(x_k,\alpha_k)$ et la commande~: 
\[
\min\ \E\sum_k\Bigl(\eta\frac{u_k^2}{\tau}+\lambda\sigma^2\tau x_k^2-\tau\,\alpha_kx_k\Bigr),\qquad x_{k+1}=x_k-u_k,\quad \alpha_{k+1}=\phi\,\alpha_k+\text{bruit},
\]
avec $x_n=0$ imposé par une forte pénalité terminale. Un problème linéaire-quadratique admet une politique optimale linéaire, $u_k=K_{k,1}x_k+K_{k,2}\alpha_k$, dont les gains se calculent par
récursion de Riccati rétrograde~; sans signal, elle reproduit exactement la trajectoire d'Almgren--Chriss. Cartea et Jaimungal (2016) ont intégré le flux d'ordres au
même cadre.

\omcode{code/firm/execcontrol/firm_execcontrol.py}{35}{47}{La récursion de Riccati de l'exécution avec un signal~: gains sur les quantités détenues et sur la dérive prévue.}\label{lst:mx:beyond-almgren-chriss:lqr}

Sur le bloc du chapitre~14 (1 million d'actions d'un titre à USD~50 sur une journée, en 78 pas de cinq minutes, $\lambda=10^{-6}$), 1\,500 journées simulées avec les mêmes
nombres aléatoires pour les deux calendriers donnent le gain présenté dans la \cref{fig:mx:beyond-almgren-chriss:signal}. Un signal de demi-vie égale à un vingtième de journée (environ
vingt minutes) et un écart-type de USD~0.5 par jour fait économiser 0.19 points de base contre 61 pour Almgren--Chriss~; un signal de demi-vie d'une journée et
un écart-type de USD~2 par jour fait économiser 41 sur 66. L'économie croît avec la force du signal et plus que linéairement avec sa persistance~: une prévision qui ne dure que quelques
pas de temps ne peut être exploitée avant de s'éteindre. Elle augmente aussi la variance du coût (écart-type de 212 points de base contre 173 pour le plus
fort signal)~: le calendrier prend du risque pour utiliser la prévision.

\begin{omfigure}
\begin{tikzpicture}
\begin{axis}[width=9cm,height=5.4cm,xmode=log,xlabel={\small demi-vie du signal (jours)},ylabel={\small économie (points de base)},
  tick label style={font=\footnotesize},xtick={0.05,0.2,1},xticklabels={0.05,0.2,1},grid=major,grid style={gray!25},ymin=0,ymax=45,
  legend columns=3,legend style={font=\footnotesize,at={(0.5,-0.3)},anchor=north,draw=none,fill=none,
  /tikz/every even column/.append style={column sep=8pt}},table/col sep=comma]
\addplot[omDef,very thick,mark=*,error bars/.cd,y dir=both,y explicit] table[x=half,y=s0,y error=e0]{figdata/microstructure/15-beyond-almgren-chriss/signal.csv};
\addplot[omThm,very thick,mark=square*,error bars/.cd,y dir=both,y explicit] table[x=half,y=s1,y error=e1]{figdata/microstructure/15-beyond-almgren-chriss/signal.csv};
\addplot[omProp,very thick,mark=triangle*,error bars/.cd,y dir=both,y explicit] table[x=half,y=s2,y error=e2]{figdata/microstructure/15-beyond-almgren-chriss/signal.csv};
\legend{USD 0.5 par jour,USD 1 par jour,USD 2 par jour}
\end{axis}
\end{tikzpicture}

\omcaption{Ce qu'une prévision de dérive fait économiser à l'exécution du bloc, en fonction de sa demi-vie, pour trois intensités (écarts-types de la dérive)~: économie moyenne
sur 1\,500 journées simulées avec nombres aléatoires communs, et son erreur type. Données~: \texttt{mx\_beyond.signal\_study}.}\label{fig:mx:beyond-almgren-chriss:signal}
\end{omfigure}

\section{Calendriers adaptatifs}

\begin{definition}[Exécution adaptative, stratégie agressive dans la monnaie]\label{def:mx:beyond-almgren-chriss:adaptive}
Une \emph{exécution adaptative}\index{exécution adaptative} (adaptive execution) modifie son calendrier en fonction de la trajectoire du cours pendant l'ordre. Une \emph{stratégie agressive dans la
monnaie}\index{stratégie agressive dans la monnaie} (aggressive-in-the-money strategy) exécute plus vite lorsque le cours a évolué dans le sens de l'ordre (à la hausse pour un vendeur) et plus lentement lorsqu'il a
évolué à rebours.
\end{definition}

Almgren--Chriss est statique~: son optimum moyenne--variance est calculé une seule fois. Lorenz et Almgren (2011) ont montré qu'une réoptimisation au fil des mouvements du cours fait mieux
au même sens moyenne--variance. L'intuition~: après un mouvement favorable, l'ordre a dégagé un gain qu'il peut dépenser sous forme d'impact pour réduire le risque restant.
La règle du chapitre recalcule à chaque pas la trajectoire d'Almgren--Chriss à partir des quantités détenues, avec une \omterm{def:mx:the-almgren-chriss-framework:urgency}{urgence} $\kappa_0e^{a\,\Delta p/(\sigma\sqrt T)}$.

\omcode{code/firm/execcontrol/firm_execcontrol.py}{78}{91}{Une stratégie agressive dans la monnaie~: la trajectoire statique recalculée à partir des quantités détenues, avec une urgence qui croît avec le gain accumulé.}\label{lst:mx:beyond-almgren-chriss:aim}

Avec $\kappa_0T=2$ et 3\,000 journées simulées, le calendrier statique coûte 60.9 points de base avec un écart-type de 90.5, un équivalent certain de 101.9. La
règle adaptative avec $a=0.25$, 0.5, 1 et 2 coûte 61.2, 62.2, 66.2 et 81.7 en moyenne, avec des écarts-types de 88.7, 86.4, 80.4 et 66.2, et des équivalents certains
de 100.5, 99.4, 98.6 et 103.6~: une adaptation modérée améliore le compromis jusqu'à 3.2~\%, une adaptation forte va trop loin.

\begin{omfigure}
\begin{tikzpicture}
\begin{axis}[width=8.5cm,height=5.2cm,xlabel={\small écart-type (points de base)},ylabel={\small coût moyen (points de base)},
  tick label style={font=\footnotesize},xmin=60,xmax=95,ymin=55,ymax=88,grid=major,grid style={gray!25},table/col sep=comma]
\addplot[omDef,very thick,mark=*] table[x=sd,y=mean]{figdata/microstructure/15-beyond-almgren-chriss/aim.csv};
\node[font=\scriptsize,anchor=north east] at (axis cs:90.5,60.9) {statique};
\node[font=\scriptsize,anchor=west] at (axis cs:80.4,66.2) {$a=1$};
\node[font=\scriptsize,anchor=west] at (axis cs:66.2,81.7) {$a=2$};
\end{axis}
\end{tikzpicture}

\omcaption{La \omterm{def:mx:beyond-almgren-chriss:adaptive}{stratégie agressive dans la monnaie} pour une adaptation croissante $a=0$ (statique), 0.25, 0.5, 1 et 2~: elle échange coût moyen contre risque moindre, avec le meilleur
équivalent certain proche de $a=1$ à cette aversion au risque. 3\,000 journées simulées. Données~: \texttt{mx\_beyond.aim\_grid}.}\label{fig:mx:beyond-almgren-chriss:aim}
\end{omfigure}

\section{Impact transitoire et bloc optimal}

\begin{definition}[Trajectoire en baignoire]\label{def:mx:beyond-almgren-chriss:bucket}
Une \emph{trajectoire en baignoire}\index{trajectoire en baignoire} (bucket-shaped trajectory) exécute un bloc au début et un autre à la fin, à rythme constant plus lent entre les deux~:
l'optimum lorsque l'impact est transitoire.
\end{definition}

Avec le noyau exponentiel d'Obizhaeva et Wang (2013, chapitre~12), une résilience $\rho=10$ par jour et 101 transactions sur la journée, le calendrier pondéré par le temps
coûte 0.0893 (en unités de l'impact instantané de l'ordre entier) et l'optimum 0.0834, soit 6.6~\% de moins~: il négocie 0.0876 de l'ordre dans chaque bloc d'extrémité,
proche de la formule $1/(\rho T+2)=0.0833$, et 0.0083 dans chaque transaction intermédiaire. Vendre l'ensemble d'un coup coûterait 0.5. L'économie est modeste parce que le
noyau décroît vite par rapport à la journée~; avec une résilience plus lente, les blocs et l'économie augmentent.

\section{Liquidité stochastique}

\begin{definition}[Liquidité stochastique]\label{def:mx:beyond-almgren-chriss:stochliq}
La \emph{liquidité stochastique}\index{liquidité stochastique} (stochastic liquidity) laisse les paramètres d'impact eux-mêmes varier au hasard pendant l'exécution, comme le volume et
la profondeur le font, de sorte que le coût de négocier à l'instant présent dépend de l'état du marché à ce même instant.
\end{definition}

Almgren (2012) a résolu l'exécution moyenne--variance lorsque la liquidité et la volatilité varient au hasard, avec des stratégies qui s'adaptent à la qualité du marché. La
version du chapitre comporte deux régimes, normal et asséché, le régime asséché ayant quatre fois l'impact temporaire~; le marché reste normal avec une probabilité de 0.9 par
pas et asséché avec 0.7~; la politique sur (quantités détenues, régime) est issue de \texttt{firm.dpsolve}, présentée dans le livre~4. À mi-journée, avec la moitié de l'ordre restant, il
vend 60\,000 actions au pas suivant lorsque le marché est normal et 20\,000 lorsqu'il est asséché. Sur 3\,000 journées simulées, il coûte 73.2 points de base contre
91.5 pour Almgren--Chriss, qui ignore les régimes~: 18.3 points de base économisés (erreur type 0.5), avec un écart-type légèrement plus faible (89.6 contre 94.6).

\section{Un canevas de théorie du contrôle}

Les quatre problèmes partagent un même canevas~: un état (les quantités détenues et ce que révèle le marché), une commande (la transaction), un coût (impact immédiat, risque et prévision
de dérive à venir), et une méthode de résolution adaptée à la structure~: une forme close pour Almgren--Chriss, une récursion de Riccati pour les signaux linéaires-quadratiques, une
inversion de matrice pour l'impact transitoire, programmation dynamique pour les régimes. Gatheral et Schied (2011) ont résolu le cas du mouvement brownien géométrique sous forme close, sous
un autre critère de risque. Toute solution renvoie un calendrier dans l'interface de \texttt{firm.acexec}, de sorte que la couche des \omterm{def:mx:the-almgren-chriss-framework:parent}{ordres enfants} ne se soucie pas de celle qui l'a produit.

\section{Tutoriel~: négocier sur une prévision}\label{tut:mx:beyond-almgren-chriss:tut}

\begin{tutorial}
\textbf{Objectif.} Ajouter un signal, une adaptation, un impact transitoire et une liquidité aléatoire à l'exécution du bloc, et chiffrer chacun.
\textbf{État final~:} les \cref{fig:mx:beyond-almgren-chriss:signal,fig:mx:beyond-almgren-chriss:aim} et les chiffres des sections 2 à 4.
\begin{tutsteps}
\item \textbf{Signal.} \texttt{firm\_execcontrol.lqr\_signal(n, tau, eta, sigma, lam, phi)}~; \texttt{simulate(policy, \dots)}~; \texttt{mx\_beyond.signal\_study()}.
\item \textbf{Adaptation.} \texttt{aim(k0, a, sigma, T, n)}~; \texttt{aim\_grid()}.
\item \textbf{Transitoire.} \texttt{transient\_study(rho)}, avec\linebreak \texttt{firm\_propagator.optimal\_liquidation} et \texttt{OWScheduler}.
\item \textbf{Liquidité.} \texttt{liquidity\_dp(X, n, tau, eta, sigma, lam, stay)} et \texttt{liquidity\_study()}~; tracer avec \texttt{fig\_beyond.py}.
\end{tutsteps}
\textbf{Pour aller plus loin.} Transmettre chaque calendrier au \texttt{ScheduleAgent} du chapitre~14 et l'exécuter dans le marché simulé sur de nombreuses sessions~; estimer
la dérive du signal avec un filtre de Kalman, au lieu de l'observer.
\end{tutorial}

\section{Construction~: contrôle de l'exécution}\label{bld:mx:beyond-almgren-chriss:execcontrol}

\begin{build}
\bfield{Rôle} Les calendriers adaptatifs qui sous-tendent les algorithmes des chapitres 16 et 28 et les politiques d'exécution apprises du livre~12, comme indices de référence dont
les optimums sont connus.
\bfield{Interface} \texttt{lqr\_signal(n, tau, eta, sigma, lam, phi)}, \texttt{simulate(policy, X, n, tau, eta, sigma, phi, s\_alpha, seed, paths)},
\texttt{OWScheduler(X, rho, T)}, \texttt{aim(k0, a, sigma, T, n)}, \texttt{liquidity\_dp(X, n, tau, eta, sigma, lam, stay)}.
\bfield{Règles} Les ventes sont positives~; le dernier pas vend ce qui reste~; les coûts se mesurent par rapport au prix d'arrivée par action~; nombres aléatoires communs entre les politiques~;
probabilités de régime en dixièmes (les nœuds de \texttt{firm.dpsolve}).
\bfield{Tests d'acceptation} \texttt{code/firm/execcontrol/tests/}~: la politique de Riccati sans signal égale la trajectoire discrète d'Almgren--Chriss~; un
signal abaisse le coût moyen~; les blocs d'Obizhaeva--Wang~; le sens de la politique adaptative et sa limite statique~; davantage de ventes lorsque le marché est liquide.
\bfield{Pour aller plus loin} Signaux filtrés par un filtre de Kalman~; volatilité stochastique~; impact transitoire avec un signal.
\end{build}

\begin{omsources}
\item P.~Lorenz and R.~Almgren, ``Mean--variance optimal adaptive execution'', \emph{Applied Mathematical Finance} 18(5), 2011.
\item J.~Gatheral and A.~Schied, ``Optimal trade execution under geometric Brownian motion in the Almgren and Chriss framework'', \emph{International
Journal of Theoretical and Applied Finance} 14(3), 2011.
\item R.~Almgren, ``Optimal trading with stochastic liquidity and volatility'', \emph{SIAM Journal on Financial Mathematics} 3(1), 2012.
\item A.~A.~Obizhaeva and J.~Wang, ``Optimal trading strategy and supply/demand dynamics'', \emph{Journal of Financial Markets} 16(1), 2013.
\item A.~Cartea and S.~Jaimungal, ``Incorporating order-flow into optimal execution'', \emph{Mathematics and Financial Economics} 10(3), 2016.
\end{omsources}

\section{Exercices}

\begin{exercise}[$\star$]\label{exo:mx:beyond-almgren-chriss:1}
Un signal a une demi-vie de 0.2~jour. Calculer sa persistance $\phi$ sur un pas de cinq minutes d'une journée de 78 pas.
\end{exercise}

\begin{exercise}[$\star$]\label{exo:mx:beyond-almgren-chriss:2}
Expliquer pourquoi le gain sur la dérive prévue est négatif dans la politique de Riccati pour un vendeur (la politique vend moins lorsque la dérive prévue est positive).
\end{exercise}

\begin{exercise}[$\star$]\label{exo:mx:beyond-almgren-chriss:3}
Avec une résilience $\rho=10$ par jour et une fenêtre d'un jour, calculer la part de l'ordre que l'optimum d'Obizhaeva--Wang vend dans chaque bloc final.
\end{exercise}

\begin{exercise}[$\star\star$]\label{exo:mx:beyond-almgren-chriss:4}
Expliquer pourquoi un signal à demi-vie courte économise si peu, même lorsqu'il est fort.
\end{exercise}

\begin{exercise}[$\star\star$]\label{exo:mx:beyond-almgren-chriss:5}
Expliquer pourquoi l'adaptation avec $a=2$ fait moins bien que le calendrier statique en équivalent certain.
\end{exercise}

\begin{exercise}[$\star\star$]\label{exo:mx:beyond-almgren-chriss:6}
Décrire ce que fait la politique de \omterm{def:mx:beyond-almgren-chriss:stochliq}{liquidité stochastique} différemment d'Almgren--Chriss, et expliquer pourquoi elle économise 18 points de base.
\end{exercise}

\begin{exercise}[$\star\star\star$]\label{exo:mx:beyond-almgren-chriss:7}
\emph{Programmation.} Lancer l'étude du signal, le signal étant observé avec bruit (la politique voit $\alpha+\epsilon$ avec $\epsilon$ d'écart-type égal à celui de
$\alpha$). Mesurer la part du gain qui subsiste, et expliquer pourquoi.
\end{exercise}

\begin{exercise}[$\star\star\star$]\label{exo:mx:beyond-almgren-chriss:8}
\emph{Trouver la faille.} «~Notre exécution avec le signal alpha a économisé 40 points de base dans le backtest, donc le signal vaut 40 points de base sur chaque ordre.~»
\end{exercise}

\section{Problème~: négocier sur une prévision}

\begin{problem}[{problème du week-end --- négocier sur une prévision}]\label{pb:mx:beyond-almgren-chriss:1}
Une équipe de recherche remet au desk d'exécution une prévision de dérive et demande ce qu'elle vaut~; le desk se demande aussi comment s'adapter aux prix et à la liquidité.

\medskip\noindent\textbf{Partie I --- Le signal.}
\begin{enumerate}
\item Écrire le problème linéaire-quadratique et son état.
\item Expliquer pourquoi la politique optimale est linéaire, et comment ses gains se calculent.
\item Vérifier que, sans signal, on retrouve Almgren--Chriss.
\item Donner l'économie pour les trois intensités aux trois demi-vies.
\end{enumerate}

\medskip\noindent\textbf{Partie II --- Adaptation.}
\begin{enumerate}[resume]
\item Définir une \omterm{def:mx:beyond-almgren-chriss:adaptive}{stratégie agressive dans la monnaie} (aggressive in the money).
\item Donner le coût, le risque et l'équivalent certain du calendrier statique.
\item Donner la règle adaptative pour $a=0.25$, 0.5, 1 et 2.
\item Expliquer pourquoi une adaptation modérée aide.
\end{enumerate}

\medskip\noindent\textbf{Partie III --- Impact transitoire et liquidité.}
\begin{enumerate}[resume]
\item Donner les coûts du calendrier pondéré dans le temps et du calendrier optimal sous le noyau transitoire, ainsi que la taille des blocs.
\item Expliquer pourquoi l'économie est faible ici.
\item Décrire les deux régimes de liquidité et les transactions de la politique dans chacun.
\item Donner le coût et l'économie de la programmation dynamique.
\end{enumerate}

\medskip\noindent\textbf{Partie IV --- Le verdict.}
\begin{enumerate}[resume]
\item Classer les quatre extensions selon l'économie réalisée ici.
\item Déterminer lesquelles exigent des données que le desk pourrait ne pas avoir.
\item Déterminer l'effet du signal sur le risque du coût.
\item Proposer une méthode pour tester la valeur du signal sur des ordres réels.
\item Énoncer le \emph{résultat nommé}~: les points de base qu'un signal alpha permet d'économiser sur un calendrier d'exécution, en fonction de la demi-vie et de l'intensité du signal.
\item Déterminer la demi-vie minimale qui rend un signal utile à intégrer dans l'exécution, avec ces chiffres.
\item Choisir l'extension à construire en premier, et justifier le choix.
\item En une phrase, résumer ce que rapporte à une exécution le fait de regarder le marché.
\end{enumerate}
\end{problem}

\section{Questions d'entretien}

\begin{interviewq}[$\star$ \iqroles{trader}]\label{iq:mx:beyond-almgren-chriss:1}
Vous vendez et votre modèle alpha indique que l'action va baisser dans l'heure qui vient. Comment votre calendrier évolue-t-il~?
\end{interviewq}

\begin{interviewq}[$\star\star$ \iqroles{chercheur}]\label{iq:mx:beyond-almgren-chriss:2}
Formuler l'exécution avec un signal d'Ornstein--Uhlenbeck comme problème de contrôle linéaire-quadratique.
\end{interviewq}

\begin{interviewq}[$\star\star$ \iqroles{chercheur}]\label{iq:mx:beyond-almgren-chriss:3}
Pourquoi l'exécution optimale sous impact transitoire a-t-elle une \omterm{def:mx:beyond-almgren-chriss:bucket}{trajectoire en baignoire}~?
\end{interviewq}

\begin{interviewq}[$\star\star$ \iqroles{trader}]\label{iq:mx:beyond-almgren-chriss:4}
Faut-il accélérer ou ralentir une exécution après que le prix a évolué en sa faveur~?
\end{interviewq}

\begin{interviewq}[$\star\star$ \iqroles{développeur}]\label{iq:mx:beyond-almgren-chriss:5}
Comment compareriez-vous équitablement deux politiques d'exécution en simulation~?
\end{interviewq}

\begin{interviewq}[$\star\star\star$ \iqroles{chercheur}]\label{iq:mx:beyond-almgren-chriss:6}
Comment estimeriez-vous la valeur, pour l'exécution, d'un signal à court horizon à partir d'ordres historiques~?
\end{interviewq}
